% PV-MEPCG / PulseVision -- equations of blueprint section 11.
% Generated by python scripts/40_export_equations.py from src/reporting/equations.py. Do not edit;
% every formula is cross-referenced to the line of code that computes it.
\documentclass{article}
\usepackage{amsmath}
\usepackage{booktabs}
\begin{document}
\section*{PV-MEPCG / PulseVision: equations and optimization formulas}
The fifteen formulas of the developer blueprint, section 11, with their symbols and the source line that implements each. PV-MEPCG / PulseVision is an academic screening and decision-support prototype, not a diagnostic tool.

\subsection*{Butterworth response}
% implemented in src/preprocessing/filters.py:318
% magnitude = 1.0 / np.sqrt(1.0 + transformed ** (2 * order))
\begin{equation}\label{eq:butterworth}
\left|H(j\omega)\right|^{2} = \frac{1}{1 + \left(\dfrac{\omega}{\omega_c}\right)^{2n}}
\end{equation}
\noindent\textit{Use:} Preprocessing.\par
\smallskip\noindent\begin{tabular}{@{}ll@{}}
\toprule Symbol & Meaning \\ \midrule
$H(j\omega)$ & frequency response of the filter \\
$\omega$ & angular frequency, in radians per second \\
$\omega_c$ & cut-off angular frequency \\
$n$ & filter order \\
\bottomrule
\end{tabular}\par
\smallskip\noindent\textit{Implemented in} \texttt{src/preprocessing/filters.py:318}.\par
\noindent\textit{Implementation note.} The equation is the low-pass prototype. The implemented filter is a band-pass, reached by the substitution W = (w\textasciicircum{}2 - w0\textasciicircum{}2) / (w BW) with w0\textasciicircum{}2 = w1 w2 and BW = w2 - w1, after bilinear prewarping of every frequency; the reference function reproduces scipy's sosfreqz of the designed filter. The filter is run forward and backward, so the applied magnitude is the square of this response.\par

\subsection*{Z-score}
% implemented in src/preprocessing/normalize.py:161
% return ((samples - samples.mean()) / std).astype(np.float32), False
\begin{equation}\label{eq:zscore}
z = \frac{x - \mu}{\sigma}
\end{equation}
\noindent\textit{Use:} Signal normalization.\par
\smallskip\noindent\begin{tabular}{@{}ll@{}}
\toprule Symbol & Meaning \\ \midrule
$x$ & a sample of the signal \\
$\mu$ & mean of the signal \\
$\sigma$ & standard deviation of the signal \\
\bottomrule
\end{tabular}\par
\smallskip\noindent\textit{Implemented in} \texttt{src/preprocessing/normalize.py:161}.\par

\subsection*{Energy}
% implemented in src/feature_extraction/time_domain.py:366
% energy = float(np.sum(signal**2))
\begin{equation}\label{eq:energy}
E = \sum_{n} x[n]^{2}
\end{equation}
\noindent\textit{Use:} Time-domain feature.\par
\smallskip\noindent\begin{tabular}{@{}ll@{}}
\toprule Symbol & Meaning \\ \midrule
$E$ & total signal energy \\
$x[n]$ & the n-th sample \\
\bottomrule
\end{tabular}\par
\smallskip\noindent\textit{Implemented in} \texttt{src/feature\_extraction/time\_domain.py:366}.\par
\noindent\textit{Implementation note.} Computed on the z-normalized signal, where it equals the number of samples; the time-domain feature time\_energy is therefore duration x sampling rate.\par

\subsection*{RMS}
% implemented in src/feature_extraction/time_domain.py:367
% rms = float(np.sqrt(energy / signal.size))
\begin{equation}\label{eq:rms}
\mathrm{RMS} = \sqrt{\frac{1}{N}\sum_{n} x[n]^{2}}
\end{equation}
\noindent\textit{Use:} Effective amplitude.\par
\smallskip\noindent\begin{tabular}{@{}ll@{}}
\toprule Symbol & Meaning \\ \midrule
$N$ & number of samples \\
$x[n]$ & the n-th sample \\
\bottomrule
\end{tabular}\par
\smallskip\noindent\textit{Implemented in} \texttt{src/feature\_extraction/time\_domain.py:367}.\par
\noindent\textit{Transcription note.} Section 11 prints the radical after the bracket rather than over it. Rendered as the root of the mean square, which is what the name, the use column and the implementation all describe.\par
\noindent\textit{Implementation note.} Computed on the z-normalized signal, where it is identically 1 (time\_rms).\par

\subsection*{Spectral centroid}
% implemented in src/feature_extraction/frequency.py:327
% centroid = librosa.feature.spectral_centroid(S=magnitude, sr=fs, n_fft=n_fft)
\begin{equation}\label{eq:spectral_centroid}
C = \frac{\sum_{k} f_k X_k}{\sum_{k} X_k}
\end{equation}
\noindent\textit{Use:} Spectral centre of mass.\par
\smallskip\noindent\begin{tabular}{@{}ll@{}}
\toprule Symbol & Meaning \\ \midrule
$C$ & spectral centroid, in hertz \\
$f_k$ & frequency of the k-th bin \\
$X_k$ & magnitude of the k-th bin \\
\bottomrule
\end{tabular}\par
\smallskip\noindent\textit{Implemented in} \texttt{src/feature\_extraction/frequency.py:327}.\par
\noindent\textit{Transcription note.} Bracketed here. The blueprint's unparenthesised form reads as a sum of ratios, which is not a centre of mass.\par
\noindent\textit{Implementation note.} Computed per short-time Fourier frame by librosa, with X\_k the frame's magnitude spectrum, and summarised as the mean and SD over frames (freq\_centroid\_mean, freq\_centroid\_std).\par

\subsection*{Spectral entropy}
% implemented in src/feature_extraction/frequency.py:166
% entropy = float(-np.sum(probabilities * np.log2(probabilities)))
\begin{equation}\label{eq:spectral_entropy}
H_s = -\sum_{k} p_k \log p_k
\end{equation}
\noindent\textit{Use:} Spectral irregularity.\par
\smallskip\noindent\begin{tabular}{@{}ll@{}}
\toprule Symbol & Meaning \\ \midrule
$H_s$ & spectral entropy \\
$p_k$ & normalized power in the k-th bin, summing to one \\
\bottomrule
\end{tabular}\par
\smallskip\noindent\textit{Implemented in} \texttt{src/feature\_extraction/frequency.py:166}.\par
\noindent\textit{Implementation note.} Implemented with base-2 logarithms on the Welch power spectrum and divided by log2 of the number of bins, so the feature lies in [0, 1] and is comparable between records whose spectra have different lengths.\par

\subsection*{Soft voting}
% implemented in src/ensemble/soft_voting.py:114
% return np.tensordot(vector / total, stack, axes=(0, 0))
\begin{equation}\label{eq:soft_voting}
P_{\text{final}}(c) = \frac{\sum_{m} w_m P_m(c)}{\sum_{m} w_m}
\end{equation}
\noindent\textit{Use:} Weighted ensemble probability.\par
\smallskip\noindent\begin{tabular}{@{}ll@{}}
\toprule Symbol & Meaning \\ \midrule
$c$ & a class \\
$m$ & an ensemble member: SVM, Random Forest or Gradient Boosting \\
$w_m$ & the weight of member m, non-negative \\
$P_m(c)$ & member m's probability for class c \\
\bottomrule
\end{tabular}\par
\smallskip\noindent\textit{Implemented in} \texttt{src/ensemble/soft\_voting.py:114}.\par
\noindent\textit{Implementation note.} Weights are normalised by their sum before averaging, so raw scores may be passed; M6 uses equal weights and M7 the weights chosen in-fold (ALG-11, ALG-12).\par

\subsection*{Prediction}
% implemented in src/ensemble/soft_voting.py:854
% return classes[np.argmax(proba, axis=1)]
\begin{equation}\label{eq:prediction}
\hat{y} = \arg\max_{c} P_{\text{final}}(c)
\end{equation}
\noindent\textit{Use:} Final class.\par
\smallskip\noindent\begin{tabular}{@{}ll@{}}
\toprule Symbol & Meaning \\ \midrule
$\hat{y}$ & the predicted class \\
$P_{\text{final}}(c)$ & the fused probability for class c \\
\bottomrule
\end{tabular}\par
\smallskip\noindent\textit{Implemented in} \texttt{src/ensemble/soft\_voting.py:854}.\par
\noindent\textit{Transcription note.} For the binary task the deployed rule is not argmax at 0.5: T50.4 selects the decision threshold inside the training fold against sensitivity and balanced accuracy. This equation describes the multiclass rule and the fixed-threshold binary reference.\par
\noindent\textit{Implementation note.} The argmax rule is the multiclass decision of SoftVotingEnsemble.predict. For the binary task the ensembles predict by a threshold chosen inside the training fold, and the deployed model predicts by argmax at the fixed 0.5 reference.\par

\subsection*{Accuracy}
% implemented in src/evaluation/metrics.py:199
% "accuracy": float((tp + tn) / (tp + tn + fp + fn)),
\begin{equation}\label{eq:accuracy}
\mathrm{Accuracy} = \frac{TP + TN}{TP + TN + FP + FN}
\end{equation}
\noindent\textit{Use:} Overall correctness.\par
\smallskip\noindent\begin{tabular}{@{}ll@{}}
\toprule Symbol & Meaning \\ \midrule
$TP$ & true positives: abnormal, predicted abnormal \\
$TN$ & true negatives: normal, predicted normal \\
$FP$ & false positives: normal, predicted abnormal \\
$FN$ & false negatives: abnormal, predicted normal \\
\bottomrule
\end{tabular}\par
\smallskip\noindent\textit{Implemented in} \texttt{src/evaluation/metrics.py:199}.\par

\subsection*{Sensitivity}
% implemented in src/evaluation/metrics.py:194
% recall_score(true, pred, pos_label=positive_label, zero_division=0)
\begin{equation}\label{eq:sensitivity}
\mathrm{Sensitivity} = \frac{TP}{TP + FN}
\end{equation}
\noindent\textit{Use:} Abnormal-case detection.\par
\smallskip\noindent\begin{tabular}{@{}ll@{}}
\toprule Symbol & Meaning \\ \midrule
$TP$ & true positives \\
$FN$ & false negatives: the abnormal cases that were missed \\
\bottomrule
\end{tabular}\par
\smallskip\noindent\textit{Implemented in} \texttt{src/evaluation/metrics.py:194}.\par

\subsection*{Specificity}
% implemented in src/evaluation/metrics.py:149
% return float(tn / denominator) if denominator else float("nan")
\begin{equation}\label{eq:specificity}
\mathrm{Specificity} = \frac{TN}{TN + FP}
\end{equation}
\noindent\textit{Use:} Normal-case detection.\par
\smallskip\noindent\begin{tabular}{@{}ll@{}}
\toprule Symbol & Meaning \\ \midrule
$TN$ & true negatives \\
$FP$ & false positives \\
\bottomrule
\end{tabular}\par
\smallskip\noindent\textit{Implemented in} \texttt{src/evaluation/metrics.py:149}.\par

\subsection*{F1-score}
% implemented in src/evaluation/metrics.py:206
% "f1": float(f1_score(true, pred, pos_label=positive_label, zero_division=0)),
\begin{equation}\label{eq:f1}
F_1 = \frac{2PR}{P + R}
\end{equation}
\noindent\textit{Use:} Precision-recall balance.\par
\smallskip\noindent\begin{tabular}{@{}ll@{}}
\toprule Symbol & Meaning \\ \midrule
$P$ & precision, $TP/(TP+FP)$ \\
$R$ & recall, which is sensitivity, $TP/(TP+FN)$ \\
\bottomrule
\end{tabular}\par
\smallskip\noindent\textit{Implemented in} \texttt{src/evaluation/metrics.py:206}.\par

\subsection*{Balanced accuracy}
% implemented in src/evaluation/metrics.py:200
% "balanced_accuracy": float(np.nanmean([sensitivity, specificity])),
\begin{equation}\label{eq:balanced_accuracy}
\mathrm{BA} = \frac{\mathrm{Sensitivity} + \mathrm{Specificity}}{2}
\end{equation}
\noindent\textit{Use:} Imbalance-aware binary metric.\par
\smallskip\noindent\begin{tabular}{@{}ll@{}}
\toprule Symbol & Meaning \\ \midrule
$\mathrm{BA}$ & balanced accuracy \\
\bottomrule
\end{tabular}\par
\smallskip\noindent\textit{Implemented in} \texttt{src/evaluation/metrics.py:200}.\par

\subsection*{Macro-F1}
% implemented in src/optimization/multi_objective.py:258
% kwargs: dict[str, Any] = {"average": "macro", "zero_division": 0}
\begin{equation}\label{eq:macro_f1}
\text{Macro-}F_1 = \frac{1}{|C|}\sum_{c \in C} F_1(c)
\end{equation}
\noindent\textit{Use:} Multiclass evaluation.\par
\smallskip\noindent\begin{tabular}{@{}ll@{}}
\toprule Symbol & Meaning \\ \midrule
$C$ & the set of classes \\
$F_1(c)$ & the F1 score of class c taken as positive \\
\bottomrule
\end{tabular}\par
\smallskip\noindent\textit{Implemented in} \texttt{src/optimization/multi\_objective.py:258}.\par
\noindent\textit{Implementation note.} Defined for the binary task too, where it averages the F1 of both classes -- not the positive-class F1 of equation 12.\par

\subsection*{Multi-objective score}
% implemented in src/optimization/multi_objective.py:323
% weights.alpha * (1.0 - float(performance))
\begin{equation}\label{eq:objective_j}
J = \alpha\left(1 - \text{Macro-}F_1\right) + \beta\left(\frac{\text{SelectedFeatures}}{138}\right) + \gamma\left(\text{NormalizedInferenceTime}\right)
\end{equation}
\noindent\textit{Use:} Optimization objective.\par
\smallskip\noindent\begin{tabular}{@{}ll@{}}
\toprule Symbol & Meaning \\ \midrule
$J$ & the cost being minimised \\
$\alpha, \beta, \gamma$ & weights on accuracy, compactness and speed \\
$\text{SelectedFeatures}$ & size of the selected subset, out of 138 \\
$\text{NormalizedInferenceTime}$ & inference time scaled against the cost model \\
\bottomrule
\end{tabular}\par
\smallskip\noindent\textit{Implemented in} \texttt{src/optimization/multi\_objective.py:323}.\par
\noindent\textit{Implementation note.} SelectedFeatures is the subset size and 138 the declared total (JWeights.n\_features\_total). NormalizedInferenceTime is the measured extraction cost of the feature families the subset still needs, divided by the cost of all families.\par

\end{document}
